\documentclass[12pt]{article}

\usepackage[utf8]{inputenc}
\usepackage[T1]{fontenc}
\usepackage[brazil]{babel}

\usepackage{geometry}
\geometry{margin=2.5cm}
\usepackage{setspace}
\onehalfspacing

\usepackage{amsmath, amssymb, amsthm}
\usepackage{mathtools}

\numberwithin{equation}{section}

\theoremstyle{plain}
\newtheorem{theorem}{Teorema}[section]
\newtheorem{proposition}[theorem]{Proposição}
\newtheorem{lemma}[theorem]{Lema}
\newtheorem{corollary}[theorem]{Corolário}

\theoremstyle{definition}
\newtheorem{definition}[theorem]{Definição}
\newtheorem{example}[theorem]{Exemplo}
\newtheorem{exercise}{Exercício}[section]

\theoremstyle{remark}
\newtheorem*{remark}{Observação}

\usepackage{hyperref}
\usepackage{csquotes}

\usepackage[
    backend=biber,
    style=authoryear,
    maxcitenames=2
]{biblatex}

\addbibresource{/mnt/c/Users/nicho/Documents/nicholasvoltani.github.io/bibliography.bib}

\newcommand{\R}{\mathbb{R}}
\newcommand{\pref}{\succeq}
\newcommand{\strictpref}{\succ}
\newcommand{\indiff}{\sim}

% ------------------------------------------------
% Title (fill manually)
% ------------------------------------------------
\title{}
\author{}
\date{}

\begin{document}

\maketitle

\hypertarget{exercuxedcio-3}{%
\section{Exercício 3}\label{exercuxedcio-3}}

!400

Seja \(x(p,w)\) a Demanda Marshalliana, tal que \[
x_{l}(p, \alpha w) = \alpha x_{l}(p,w)
\] (homogênea de grau 1 no tocante à Renda, para dado bem \(l\)).

Derivando com relação a \(\alpha\) e tomando \(\alpha=1\) ao final,
temos \$\$ \begin{align*}
\frac{ \partial x_{l} }{ \partial w } w &= x_{l} \\
\therefore \frac{ \partial x_{l} }{ \partial w } \frac{w}{x_{l}} \equiv \epsilon_{lw} &= 1

\end{align*} \$\$ onde \(\epsilon_{lw}\) é a Elasticidade-Renda da
Demanda.

Vetorialmente, temos{[}\^{}1{]} \[
\nabla_{w} x = \frac{x}{w}
\] Sendo \(w\) dado (assim como \(p\)), tomando
\(\alpha = \frac{1}{w}\), temos \[
\nabla_{w} x = \frac{1}{w} x(p, 1)
\] Ou seja, dados \((p, w)\), temos que a curva \(x\) no tocante à renda
--- i.e.~Curva de Engel --- tem inclinação constante,
\textbf{é uma reta}, de coeficiente \(\frac{x(p,1)}{w}\).

\hypertarget{exemplos}{%
\subsubsection{Exemplos}\label{exemplos}}

\begin{enumerate}
\def\labelenumi{\arabic{enumi})}
\tightlist
\item
  Substitutos Perfeitos Caso \(p_{2}>p_{1}\), tem-se que se substitui
  totalmente o consumo de \(2\) por \(1\). Portanto, pela restrição
  orçamentária, temos \[
  x_{1} = \frac{w}{p_{1}}
  \]
\end{enumerate}

Logo, \(\frac{ \partial x_{1} }{ \partial w } = \frac{1}{p_{1}}\) é a
inclinação da curva de Engel de bens substitutos perfeitos.

\begin{enumerate}
\def\labelenumi{\arabic{enumi})}
\setcounter{enumi}{1}
\item
  Complementares Perfeitos Supondo-se bens complementares \(1:1\), temos
  que \(x_{1}=x_{2}\equiv x\). Pela restrição orçamentária, teremos \[
  x = \frac{w}{p_{1}+p_{2}} = x_{1} = x_{2}
  \] Logo,
  \(\frac{ \partial x }{ \partial w } = \frac{1}{p_{1}+p_{2}}\).
\item
  Função de Cobb-Douglas Tendo bens com coeficientes \(\alpha_{i}\) tais
  que \(\sum_{i} \alpha_{i} = 1\), temos que suas demandas serão \[
  x_{i} = \frac{\alpha_{i}}{p_{i}} w 
  \] com inclinações
  \(\frac{ \partial x_{i} }{ \partial w } = \frac{\alpha_{i}}{p_{i}}\).
\end{enumerate}

\hypertarget{exercuxedcio-5}{%
\section{Exercício 5}\label{exercuxedcio-5}}

!500

Dados os bens \(x\) e \(h\){[}\^{}2{]}, onde \(x\) possui preço \(p\) e
o tempo de trabalho possui salário \(s\), temos que a restrição
orçamentária é de que \[
\underbrace{ p\cdot x }_{ \text{Consumo de bens} } \leq \underbrace{ (24-h)\cdot s }_{ \text{Salário} }
\] o que é equivalente a \[
\underbrace{ p \cdot x + s \cdot h }_{ \text{Custos (+ de oportunidade)} } \leq \underbrace{ 24 \cdot s }_{ \text{Salário máximo} }
\]

!500 Fonte: VARIAN, p.~175.

\printbibliography

\end{document}
